\section{Proposed Work / System}

\subsection{Architectural Vision of SIGS-RS}
To resolve the vulnerabilities inherent in legacy municipal telemetry pipelines, this project presents the **Secure IoT Gateway Simulation for Resource Sensing (SIGS-RS)**. SIGS-RS is an asynchronous, event-driven edge gateway and security middleware engine specifically designed to process high-volume sensor streams across water, electrical power, and air quality infrastructure while enforcing non-negotiable security, privacy, and reliability guarantees.

Unlike monolithic ingestion engines that couple transport ingestion directly to database transactions, SIGS-RS enforces an asynchronous, multi-tiered architecture that separates ingestion transport, in-memory buffering, two-stage sequential security middleware, persistence protection, and management controls into cleanly isolated subsystems.

\subsection{The Fundamental Security Invariant}
The cornerstone of the SIGS-RS architecture is the **Sequential Processing Invariant**:
\begin{equation}
\label{eq:invariant}
\text{Raw Ingress} \longrightarrow \text{IngestionQueue} \longrightarrow \text{Inspector (Stage 1)} \longrightarrow \text{Anonymizer (Stage 2)} \longrightarrow \text{Circuit Breaker} \longrightarrow \text{Persistence}
\end{equation}

This processing pipeline establishes four strict operational rules:
\begin{enumerate}[label=(\roman*)]
    \item \textbf{Zero Raw Persistence:} Under no operational condition does uninspected or unsanitized telemetry enter persistent database storage. All ingress data must traverse the middleware pipeline.
    \item \textbf{Stage Ordering and Pre-Drop Invariant:} The Anomalous Packet Inspector strictly precedes the Privacy Anonymizer. If an incoming packet exceeds the configured sliding-window rate limit (50 packets / 10~s per sensor) or contains thermodynamically impossible values (e.g., negative flow or out-of-bounds AQI), the packet is \emph{immediately discarded}. It is never passed to the Privacy Anonymizer, eliminating wasted computational cycles and ensuring forensic anomaly logs record the original, unredacted packet identifiers.
    \item \textbf{Fail-Secure Anonymization:} Verified packets traversing the Privacy Anonymizer have all spatial coordinates (\texttt{latitude}, \texttt{longitude}, \texttt{geo}) and hardware identifiers (\texttt{mac\_address}, \texttt{device\_id}) permanently stripped. If the anonymization routine encounters an unhandled exception or detects residual PII fields, the entire packet is discarded fail-securely.
    \item \textbf{Persistence Protection:} Sanitized packets can only be committed to the database if the persistence Circuit Breaker is in the \texttt{CLOSED} or \texttt{HALF-OPEN} state. If database write latency exceeds 2,000~ms, the circuit breaker trips to \texttt{OPEN}, and the gateway immediately suspends telemetry ingress to protect gateway memory and host resources.
\end{enumerate}

\subsection{Core Engineering Innovations}

\subsubsection{Asynchronous Queue Buffering and Backpressure Handling}
Incoming telemetry from MQTT (TLS on port 8883) and HTTP REST (\texttt{POST /api/v1/telemetry}) is accepted by non-blocking ingress handlers that validate content length ($\le 512$~KB per SRS REQ-TI-3) and JSON syntax before placing packets into an in-memory \texttt{IngestionQueue} (SRS REQ-TI-4). The queue decouples the high-velocity arrival rate of virtual sensors from the batch-processing rate of the database worker, absorbing transient traffic spikes without stalling network sockets.

\subsubsection{Two-Tier Sliding-Window Implementation}
To guarantee rigorous rate-limiting across diverse environments, SIGS-RS supports two interchangeable sliding-window engines:
\begin{itemize}
    \item \textbf{In-Memory Deque Engine:} Maintains an in-process dictionary of double-ended queues (\texttt{collections.deque}) tracking packet timestamps for each sensor. Timestamps outside the sliding window $[t - 10\text{s}, t]$ are evicted in $O(1)$ amortized time, providing microsecond evaluation latency ($p99 < 0.02$~ms) without external dependencies.
    \item \textbf{Distributed Redis Lua Engine:} Utilizes Redis Sorted Sets (\texttt{ZSET}) executed via an atomic Lua script for horizontally scaled deployments. This guarantees that multiple gateway pods share identical rate-limiting state without race conditions.
\end{itemize}

\subsubsection{Automated Regulatory Compliance Lockout (SRS BR-3)}
When privacy compliance guidelines are revised through the administrative API (\texttt{PUT /api/v1/config/compliance/guideline}), SIGS-RS automatically activates an immediate compliance lockout. The gateway transitions into a fail-closed state, refusing all incoming telemetry (HTTP 503) across all instances. Testing remains suspended until an authorized Compliance Auditor or Security Engineer posts a formal audit signoff (\texttt{POST /api/v1/config/compliance/audit-signoff}), guaranteeing end-to-end regulatory oversight.

\subsubsection{Database-Enforced Audit Immutability (SRS REQ-SE-4)}
To satisfy legal compliance mandates, all security-relevant occurrences (system startup, authentication failures, rate drops, anomaly drops, and compliance lockouts) are committed to \texttt{audit.security\_log}. Immutability is enforced at the database engine level via the PostgreSQL trigger function \texttt{audit.prevent\_tampering()}, which unconditionally blocks all \texttt{UPDATE} and \texttt{DELETE} statements.

\subsection{Verification Boundaries and Technical Disclosures}
In accordance with academic rigor, the proposed architecture clearly delineates between features verified in the local test environment and those requiring production infrastructure:
\begin{itemize}
    \item \textbf{Local Verification Baseline:} The complete pipeline, sliding-window algorithms, physical validators, PII stripping, circuit breaker transitions, and database immutability triggers are 100\% implemented and verified locally using native Python 3.12, \texttt{InMemoryDatabasePool}, and automated pytest suites.
    \item \textbf{Production Cluster Scope:} Achieving sustained 50,000 packets/second across 10,000 sensors (REQ-PF-1), hardware-level WORM storage, and enterprise OAuth 2.0 Identity Provider federation are designed as production cluster targets (requiring 16 vCPUs, 32 GB RAM, and managed Redis/TimescaleDB clusters per SRS §3.2) and are truthfully documented as unverified on the local host.
\end{itemize}
